\section{Related Work}
\textbf{Code as the reasoning medium.} PAL~\cite{gao2022pal} and Program of Thoughts~\cite{chen2022pot} have the model write a program and leave the arithmetic to an interpreter. ToRA~\cite{gou2023tora} combines natural-language reasoning with program tools for mathematical problems. We share the view that the model should not do the arithmetic, but we limit the program to calls into a fixed set of functions.

\textbf{Data-science agents.} Data Interpreter~\cite{hong2024datainterpreter} is the closest work and our base paper. It models a project as a graph of tasks, each task as a graph of code actions, refines the task graph when execution fails, and retrieves tools by their schemas before writing code that combines them. The paper reports 94.93\% on InfiAgent-DABench~\cite{hu2024dabench} with gpt-4o, and notes that its experiments cover entry-level Kaggle data and that running generated code carries a security risk. DS-Agent~\cite{guo2024dsagent} applies case-based reasoning to automate data-science tasks, and AutoGen~\cite{wu2023autogen} provides a framework for multi-agent conversations that can execute code.

\textbf{Tool use and planning.} Toolformer~\cite{schick2023toolformer} trains a model to decide when to call external tools, and HuggingGPT~\cite{shen2023hugginggpt} uses a language model as a controller that plans tasks and calls other models. ReAct~\cite{yao2022react} interleaves reasoning steps with actions, and Reflexion~\cite{shinn2023reflexion} lets an agent use verbal feedback from earlier attempts. Our control loop of planning, executing, inspecting and replanning follows the same pattern, with deterministic tools in the action step.

\textbf{Position.} Compared with Data Interpreter, our design differs in four ways: the action space is a registry of eight tested functions instead of generated code; every reported number carries an evidence record that can be checked; the evaluation uses an independent oracle; and the planner is a 4B-parameter model that runs locally. We compare designs, not scores. The tasks, the models and the benchmarks are different, and no number in this paper is meant to be set against the figures in~\cite{hong2024datainterpreter}.
